\documentclass[10pt, titlepage]{article}

\usepackage[margin=1in]{geometry}
\usepackage{fancyhdr}
\usepackage{amsfonts}
\usepackage{amsmath}
\usepackage{amssymb}
\usepackage{booktabs}
%\usepackage{tikz}
\usepackage{float}
\usepackage{bm}
\usepackage{pgfplots}


\title{\normalsize PLUS Session Worksheet \\[1em] \Large \textit{MATH 1554:} Linear Algebra}
\author{\textit{Covering Sections:} 2.2 \& 2.3}
\date{Fall 2023}

% ------- Set up header -------
\renewcommand{\sectionmark}[1]{\markright{#1}} % from https://tex.stackexchange.com/a/317393/227665
\pagestyle{fancy}
\fancyhead[L]{A. M., Fall 2023}
%\fancyhead[C]{\textit{\nouppercase{\rightmark}}}
\fancyhead[C]{\textit{Inverses (2.2 \& 2.3)}}
\fancyhead[R]{\textbf{\thepage}}
% -------------------------------------------------------

% ------- Set up question environment -------
\newcounter{question}
\setcounter{question}{1}
\newenvironment{newquestion}
{
	\noindent\textbf{Question \thequestion}
}
{
\stepcounter{question}
\par
}
% -------------------------------------------------------

% ------- Set up Main Idea box environment -------
\newcounter{mainideas}
\setcounter{mainideas}{1}
\newcommand{\mainidea}{\textbf{Main Idea \#\themainideas \stepcounter{mainideas}}}
\setlength\fboxsep{1em}
%\newcommand{\midea}[2]{\begin{center}\fbox{\parbox{#1}{\centering \mainidea #2}} \end{center}} % Could be environment but I couldn't bother

% The way this works is that the savebox collects whatever was between the minipage and the end of it, then in the last line of the definition of the environment
% we regurgigate whatever was collected into an fbox which we center.
\newsavebox{\mybox}
\newenvironment{mideea}[1]
{%
\begin{lrbox}{\mybox}%
	\begin{minipage}{#1\textwidth}
		{\centering \mainidea \par}

}
{\end{minipage}\end{lrbox}\begin{center}\fbox{\usebox{\mybox}}\end{center}}
% -------------------------------------------------------


% ------- Set up math stuff -------
\renewcommand{\vec}[1]{\bm{#1}}
\DeclareMathOperator{\vspan}{Span}
\DeclareMathOperator{\row}{row}
\DeclareMathOperator{\col}{col}
\DeclareMathOperator{\colspace}{Col}
\DeclareMathOperator{\rowspace}{Row}
\DeclareMathOperator{\nullspace}{Nul}
\DeclareMathOperator{\rank}{rank}
\DeclareMathOperator{\proj}{proj}
% -------------------------------------------------------

% ------- Set up circle -------
\newcommand{\bubble}{\tikz[baseline=-0.625ex]\draw[radius=0.5em, fill=white] (0,0) circle;}
\newcommand{\dbubble}{\tikz[baseline=-0.625ex]\draw[radius=0.5em, fill=black] (0,0) circle;}
% -------------------------------------------------------

\makeatletter
\setlength{\@fptop}{0pt}
\makeatother

\begin{document}
\maketitle

\section{\centering T/F Questions}
\begin{newquestion}
	For each statement, mark the statement \emph{true} if it's \emph{always} true, or \emph{false} otherwise.
	\begin{table}[!h]
		\centering
		\begin{tabular}{c@{\hspace{2.5em}}cp{0.75\textwidth}}
			true & false & \\
			\specialrule{\lightrulewidth}{0.25em}{1.25em}
			\bubble & \bubble & If \(A\) is an \(n\times n\) matrix and \(AB = I_n\) for some \(n \times n\) matrix \(B\), then \(A\) is invertible. \\[0.75em]
			\bubble & \bubble & If \(B\) is row equivalent to an invertible \(n \times n\) matrix \(A\), then \(B\) is invertible.\\[0.75em]
			\bubble & \bubble & If \(A\) is an invertible \(n \times n\) matrix, then for square matrices \(B, C\), \(AB = AC\) implies \(B = C\).\\[0.75em]
			\bubble & \bubble & If the columns of a \(n \times n\) matrix \(A\) span \(\mathbb{R}^n\), then \(A\vec{x} = \vec{0}\) only has the trivial
						      solution.\\[0.75em]
			\bubble & \bubble & If an \(n \times n\) matrix \(A\) has the property that for any vectors \(\vec{x}, \vec{y} \in \mathbb{R}^n\) that are not the
						      same, \(A\vec{x} \ne A\vec{y}\), then the matrix is invertible.\\[0.75em]
			\bottomrule
		\end{tabular}
	\end{table}
\end{newquestion}

\section{\centering Short Answer Questions}
\begin{newquestion}
	Answer each part so that you may be able to come up with an answer for the following question.
	\begin{enumerate}
		\item Is the statement \[\text{``For any matrix} \ A,B,C \in \mathbb{R}^{n \times n}, \ \text{if} \ AB = CA = I_n \ \text{then} \ B = C,"\] always true?
			\begin{enumerate}
				\item Note that \(AB = I_n\). How might you be able to use the information provided to get to \(B = C\)? (You can also start from 
					\(CA = I_n\) and reach \(B = C\) as well.)
				\item What does this hint to us about the of inverses of a square matrix?
			\end{enumerate}

		\item How does the standard matrix of the inverse of a linear transformation from \(\mathbb{R}^n\) to \(\mathbb{R}^n\) look like?
			\begin{enumerate}
				\item What is the standard matrix of the transformation that rotates every vector in \(\mathbb{R}^3\) counterclockwise by \(\pi/2\) radians
					around the \(x_1\)-axis?
				\item Keeping in mind what it means to be the inverse of a transformation, what would the standard matrix of the inverse of the above
					transformation be?
			\end{enumerate}

		\item How do elementary matrices relate to inverses of matrices?
			\begin{enumerate}
				\item Consider the matrix \[A = \begin{bmatrix}1 & 0 & -6 \\ 1 & 1 & -6 \\ 1 & 1 & -5\end{bmatrix}.\] Here is the row-reduction process:
					\[\begin{bmatrix}1 & 0 & -6 \\ 1 & 1 & -6 \\ 1 & 1 & -5\end{bmatrix} \stackrel{E_1}{\sim} \begin{bmatrix}1 & 0 & -6 \\ 1 & 1 & -6 \\ 0 & 0 & 1\end{bmatrix} \stackrel{E_2}{\sim} \begin{bmatrix}1 & 0 & -6 \\ 0 & 1 & 0 \\ 0 & 0 & 1\end{bmatrix} \stackrel{E_3}{\sim} \begin{bmatrix}1 & 0 & 0 \\ 0 & 1 & 0 \\ 0 & 0 & 1\end{bmatrix}.\] Find  the elementary matrices \(E_1, E_2, E_3\), and then take their products. Then, find \((E_3E_2E_1)A\).

				\item Now, what can you say about the connection between compositions of elementary matrices and inverses of invertible matrices?
			\end{enumerate}
		
		\item[\textit{Extra}] If \(A = CB\), with \(C\) being invertible and \(A, B, C \in \mathbb{R}^{n \times n}\), explain how \(A\) and \(B\) are 
					     row-equivalent.
	\end{enumerate}
\end{newquestion}

\fancyhead[C]{\textit{Solutions}}

\section{\centering Solutions}
\subsection{Solutions to T/F Questions}
	\begin{table}[!h]
		\centering
		\begin{tabular}{c@{\hspace{2.5em}}cp{0.75\textwidth}}
			true & false & \\
			\specialrule{\lightrulewidth}{0.25em}{1.25em}
			\dbubble & \bubble & If \(A\) is an \(n\times n\) matrix and \(AB = I_n\) for some \(n \times n\) matrix \(B\), then \(A\) is invertible. \\[0.75em]
			\dbubble & \bubble & If \(B\) is row equivalent to an invertible \(n \times n\) matrix \(A\), then \(B\) is invertible.\\[0.75em]
			\dbubble & \bubble & If \(A\) is an invertible \(n \times n\) matrix, then for square matrices \(B, C\), \(AB = AC\) implies \(B = C\).\\[0.75em]
			\dbubble & \bubble & If the columns of a \(3 \times 3\) matrix \(A\) span \(\mathbb{R}^n\), then \(A\vec{x} = \vec{0}\) only has the trivial
						       solution.\\[0.75em]
			\dbubble & \bubble & If an \(n \times n\) matrix \(A\) has the property that for any vectors \(\vec{x}, \vec{y} \in \mathbb{R}^n\) that are not the
						      same, \(A\vec{x} \ne A\vec{y}\), then the matrix is invertible.\\[0.75em]
			\bottomrule
		\end{tabular}
	\end{table}

\subsection{Solutions to Short Answer Questions}
\begin{enumerate}
	\item The statement is true!
		\begin{enumerate}
			\item Since \(AB = I_n\), \(CAB = CI_n = C\). We know that \(CA = I_n\), so the equation \(CAB = CI_n\) reduces to \(B = C\). If you started off
				with \(CA = I_n\), then you can follow a similar argument to show that \(B = C\) as well.
			\item This tells you that the left and right inverses for square matrices are the same. The question from the \textbf{T/F} section also hints to the fact
				that inverse of the matrix is unique.
		\end{enumerate}
	\item If \(T \colon \mathbb{R}^n \to \mathbb{R}^n\) is defined using the matrix \(T(\vec{x}) = A\vec{x}\), where \(A \in \mathbb{R}^{n \times n}\) is an 
		invertible matrix, then the inverse of \(T\), \(T^{-1}\), is defined by \(T^{-1}(\vec{x}) = A^{-1}\vec{x}\).
		\begin{enumerate}
			\item It is \[\begin{bmatrix}1 & 0 & 0\\ 0 & 0 & -1\\ 0 & 1 & 0\end{bmatrix}.\]
			\item The standard matrix of the inverse would then be \(\begin{bmatrix}1 & 0 & 0\\ 0 & 0 & -1\\ 0 & 1 & 0\end{bmatrix}^{-1}\). Geometrically,
				this standard matrix should rotate vectors in \(\mathbb{R}^3\) by \(\pi/2\) radians \emph{clockwise} around the \(x_1\)-axis. The matrix that
				does this is \[\begin{bmatrix}1 & 0 & 0\\ 0 & 0 & 1\\ 0 & -1 & 0\end{bmatrix}.\]
		\end{enumerate}
	\item The product of a sequence of \(n \times n\) elementary matrices defines an inverse for some matrix \(A \in \mathbb{R}^{n \times n}\).
		\begin{enumerate}
			\item \[E_1 = \begin{bmatrix}1 & 0 & 0 \\ 0 & 1 & 0 \\ 0 & -1 & 1\end{bmatrix}, \quad E_2 = \begin{bmatrix}1 & 0 & 0 \\ -1 & 1 & 0 \\ 0 & 0 & 1\end{bmatrix}, \quad E_3 = \begin{bmatrix}1 & 0 & 6 \\ 0 & 1 & 0 \\ 0 & 0 & 1\end{bmatrix}.\] \[E_3E_2E_1 = \begin{bmatrix}1 & -6 & 6 \\ -1 & 1 & 0 \\ 0 & -1 & 1\end{bmatrix},\quad (E_3E_2E_1)A = I_3.\]
			\item We see that \(E_3E_2E_1 A = I_n \implies (E_3E_2E_1) I = A^{-1}\). The elemntary matrices that reduce \(A\) to the identity matrix are the 
				same to those that reduce \(I_3\) to \(A^{-1}\).
		\end{enumerate}
	\item[\textit{Extra}] If \(C\) is invertible, then it is a product of elementary matrices, so we're applying row operations on \(B\) to get \(A\). Hence, \(A\) is row
				     equivalent to \(B\).
\end{enumerate}

\end{document}